\documentclass[11pt,a4paper]{article}

% ── Packages ──────────────────────────────────────────────────────────────────
\usepackage[margin=2.5cm]{geometry}
\usepackage{amsmath, amssymb}
\usepackage{booktabs}
\usepackage{longtable}
\usepackage{array}
\usepackage{xcolor}
\usepackage{listings}
\usepackage{hyperref}
\usepackage{enumitem}
\usepackage{titlesec}
\usepackage{parskip}
\usepackage{microtype}
\usepackage{fancyhdr}
\usepackage{graphicx}
\usepackage{multirow}

% ── Colours ───────────────────────────────────────────────────────────────────
\definecolor{commentgreen}{RGB}{34,139,34}
\definecolor{keywordblue}{RGB}{0,0,200}
\definecolor{stringred}{RGB}{163,21,21}
\definecolor{lightgray}{RGB}{245,245,245}
\definecolor{accentblue}{RGB}{30,90,160}
\definecolor{passgreen}{RGB}{0,128,0}
\definecolor{failred}{RGB}{180,0,0}
\definecolor{weakorange}{RGB}{180,100,0}

% ── Code listing style ────────────────────────────────────────────────────────
\lstset{
  basicstyle=\ttfamily\small,
  keywordstyle=\color{keywordblue}\bfseries,
  commentstyle=\color{commentgreen},
  stringstyle=\color{stringred},
  backgroundcolor=\color{lightgray},
  frame=single,
  framesep=4pt,
  rulecolor=\color{gray!40},
  breaklines=true,
  showstringspaces=false,
  tabsize=2,
  captionpos=b,
}

% ── Section formatting ────────────────────────────────────────────────────────
\titleformat{\section}{\large\bfseries\color{accentblue}}{}{0em}{}[\titlerule]
\titleformat{\subsection}{\normalsize\bfseries}{}{0em}{}
\titleformat{\subsubsection}{\normalsize\itshape}{}{0em}{}

% ── Header / footer ───────────────────────────────────────────────────────────
\pagestyle{fancy}
\fancyhf{}
\rhead{\small N3 DVOL Quant Framework}
\lhead{\small Project Description}
\rfoot{\small\thepage}

% ── Hyperref ──────────────────────────────────────────────────────────────────
\hypersetup{
  colorlinks=true,
  linkcolor=accentblue,
  urlcolor=accentblue,
}

% ── Custom commands ───────────────────────────────────────────────────────────
\newcommand{\code}[1]{\texttt{#1}}
\newcommand{\Pass}{\textcolor{passgreen}{\textbf{Pass}}}
\newcommand{\Fail}{\textcolor{failred}{\textbf{Fail}}}
\newcommand{\Weak}{\textcolor{weakorange}{\textbf{Weak}}}
\newcommand{\Survivor}{\textcolor{passgreen}{\textbf{Survivor}}}
\newcommand{\Killed}{\textcolor{failred}{\textbf{Killed}}}
\newcommand{\comment}[1]{%
  \begin{quote}
    \small\itshape\textbf{Commentary:} #1
  \end{quote}
}

% ─────────────────────────────────────────────────────────────────────────────

\begin{document}

\begin{titlepage}
  \centering
  \vspace*{3cm}
  {\huge\bfseries N3 DVOL Fear Resolution}\\[0.5em]
  {\large Quantitative Research \& Paper Trading System}\\[2em]
  {\large\itshape Complete Project Description \& Commentary}\\[3em]
  \rule{0.5\textwidth}{0.4pt}\\[2em]
  {\normalsize BTCUSDT USDT-M Perpetual Futures --- Binance / Deribit}\\[0.5em]
  {\normalsize As of 2026-05-13}\\[3em]
  \vfill
  {\small\textit{Research and paper trading only. Not financial advice.}}
\end{titlepage}

\tableofcontents
\newpage

% =============================================================================
\section{Research Pipeline}
% =============================================================================

\subsection{Overview and Methodology}

This project implements a full quantitative signal research pipeline for BTCUSDT USDT-M perpetual futures on Binance. The workflow follows a strict \textbf{kill-first methodology}: every candidate signal is subjected to increasingly demanding statistical tests and is killed the moment it fails. Only signals that survive every gate are promoted to paper trading.

This approach directly counters the most common failure mode in systematic trading research --- overfitting through survivorship bias. The pipeline is:

\begin{center}
\small
\texttt{Hypothesis} $\rightarrow$ \texttt{IC screen} $\rightarrow$ \texttt{Non-overlapping OOS} $\rightarrow$ \texttt{Kill attempt (extended OOS)}
$\rightarrow$ \texttt{Regime analysis} $\rightarrow$ \texttt{Position backtest} $\rightarrow$ \texttt{Paper trading}
\end{center}

\subsection{Data}

Three primary datasets, all stored as Apache Parquet:

\begin{longtable}{@{}llll@{}}
\toprule
\textbf{File} & \textbf{Contents} & \textbf{Range} & \textbf{Rows} \\
\midrule
\endhead
\code{BTCUSDT\_1m\_klines.parquet} & 1-minute OHLCV & 2023-01-01 -- 2026-05-13 & 1,768,728 \\
\code{BTC\_deribit\_dvol\_1h.parquet} & Deribit DVOL (hourly) & 2022-01-01 -- 2026-05-13 & 38,233 \\
\code{BTCUSDT\_funding.parquet} & Binance 8h funding & 2023-01-01 -- 2026-05-13 & 3,685 \\
\code{BTCUSDT\_oi\_5m.parquet} & Open interest (5m) & 2022-01-01 -- 2026-03-29 & --- \\
\code{BTCUSDT\_bybit\_funding.parquet} & Bybit funding (cross-check) & --- & --- \\
\bottomrule
\end{longtable}

The 1-minute granularity is essential: 24-hour forward returns are computed by shifting 1440 bars, and funding PnL calculation requires exact intraday timestamps to correctly attribute which 8-hour settlements fall within a given hold window. Using Deribit DVOL data from 2022 (a severe bear market with extreme IV spikes) gives the 30-day rolling z-score genuine historical range --- without it, the normalisation would be calibrated only on moderate-IV environments and the thresholds would shift.

\comment{The dataset construction is sound. Using Parquet over CSV is the right call at this scale --- Parquet's columnar compression reduces the 1.77M-row klines file to a fraction of the equivalent CSV size while being faster to read. The extension to 2026 via \code{data/extend\_to\_2026.py} is important: stopping at 2024 would have masked the DVOL suppression regime of 2025-H2 (mean DVOL = 41.6), which is one of the most informative out-of-sample periods.}

\subsection{Train / Validation / Test Split}

\begin{center}
\begin{tabular}{@{}ll@{}}
\toprule
\textbf{Period} & \textbf{Dates} \\
\midrule
Train & 2023-01-01 to 2023-12-31 \\
Validation & 2024-01-01 to 2024-06-30 \\
Test & 2024-07-01 to 2024-12-31 \\
Walk-forward extension & through 2026-05-13 \\
\bottomrule
\end{tabular}
\end{center}

\comment{This is a correct structure. The train set is used only for parameter intuition; the frozen rule is evaluated exclusively on data after 2024-01-01. The critical detail is \emph{non-overlapping} IC computation --- using every 1440th bar rather than every bar. Overlapping daily returns on 1-minute data inflate the effective sample size by $\approx 1440\times$, producing false statistical significance. The non-overlapping IC of $+0.112$ versus the naive overlapping $+0.089$ is a real and important diagnostic.}

% ─────────────────────────────────────────────────────────────────────────────
\subsection{Phase 1 --- Signal Kill Log (Nine Signals Killed)}

\begin{longtable}{@{}lllp{6.5cm}@{}}
\toprule
\textbf{Script} & \textbf{Signal} & \textbf{Hypothesis} & \textbf{Kill reason} \\
\midrule
\endhead
07 & A/B & OFI momentum/reversion & IC $\approx -0.01$ to $-0.04$; positive at maker cost but not taker \\[4pt]
09 & C & Funding premium reversion & Regime-sensitive; destroyed by 2024 bull market \\[4pt]
08 & E (VPIN) & Order toxicity asymmetry & No robust edge after adjusting for look-ahead \\[4pt]
10 & F & Basis spread divergence & Spurious --- artefact of funding construction \\[4pt]
11 & G & Premium index signal & Below breakeven consistently \\[4pt]
12--15 & H1/H2 & Funding 1h/8h mean-reversion & Survives maker cost at 1h in-sample; destroyed OOS 2024 \\[4pt]
09 & I & OI divergence & IC below noise threshold across all lookbacks \\
\bottomrule
\end{longtable}

\comment{The H1/H2 kill story is the most instructive. H1 showed a 1-hour IC that survived maker-cost breakeven analysis in-sample. Five dedicated scripts (12--15) were written before it was killed. That is the correct approach --- most researchers would have promoted H1 after the first positive backtest. The 2024 bull run made funding mean-reversion unviable: persistent positive funding became the norm rather than the exception, and long bias dominated the return distribution.}

% ─────────────────────────────────────────────────────────────────────────────
\subsection{Phase 2 --- Signal Screen (Script 16)}

Nine new signals screened simultaneously using rolling IC ratio as a quick filter, followed by block bootstrap significance testing on the survivors:

\begin{center}
\begin{tabular}{@{}llll@{}}
\toprule
\textbf{Signal} & \textbf{IC ratio} & \textbf{Bootstrap $p$} & \textbf{Outcome} \\
\midrule
N3 (DVOL 30d z-score) & 3.89$\times$ & 0.000 & \Survivor \\
N2 (DVOL 8h slope) & 2.37$\times$ & 0.000 & \Killed\ at daily bootstrap \\
N1 (DVOL level) & 1.20$\times$ & 0.000 & \Killed\ at daily bootstrap \\
D4 (funding $\times\ \Delta$OI) & 0.34$\times$ (real) & 0.634 & \Killed\ (construction artefact) \\
D1, D2, K, L1, L2 & $< 0.13\times$ & all fail & \Killed\ immediately \\
\bottomrule
\end{tabular}
\end{center}

\comment{The N2/N3 separation is subtle and important. N2 (8-hour DVOL slope) and N3 (30-day z-score) are both DVOL-based but measure fundamentally different things. N2 measures the direction of recent change; N3 measures how far current IV deviates from its recent history. N2 failing at the daily bootstrap while N3 passes means the edge is in the \emph{relative level} of IV, not in its recent trend. The D4 kill is also notable --- the funding $\times\ \Delta$OI product initially showed promise but was an artefact of how the feature was constructed, not a real signal.}

% ─────────────────────────────────────────────────────────────────────────────
\subsection{Signal N3 Definition}

The N3 signal is the 30-day rolling z-score of the Deribit BTC implied volatility index (DVOL):

\begin{equation}
z_t^{N3} = \frac{\mathrm{DVOL}_t - \mu_{30\mathrm{d}}}{\sigma_{30\mathrm{d}}}
\end{equation}

where $\mu_{30\mathrm{d}}$ and $\sigma_{30\mathrm{d}}$ are computed over a rolling window of $30 \times 24 = 720$ hourly DVOL bars.

The signal is \textbf{long only} and \textbf{positive}: a high z-score predicts positive 24-hour BTC returns. The proposed mechanism is \emph{fear resolution} --- elevated implied volatility reflects market stress, and as fear subsides over the subsequent 24 hours, prices tend to normalise upward. This is a distinct mechanism from funding carry, order flow imbalance, or basis arbitrage.

\comment{The 30-day lookback is not arbitrary: it is long enough to capture a meaningful baseline while staying within one market regime. The z-score formulation (rather than the raw level or slope) is the correct form for a mean-reversion signal --- it expresses the current state relative to recent history, which is what the fear-resolution mechanism actually predicts.}

% ─────────────────────────────────────────────────────────────────────────────
\subsection{N3 Rigorous Validation (Script 18)}

\begin{center}
\begin{tabular}{@{}ll@{}}
\toprule
\textbf{Metric} & \textbf{Value} \\
\midrule
Non-overlapping OOS IC (24h) & $+0.112$ \\
Block bootstrap $p$-value & $0.039$ \\
OOS IC (48h) & $+0.158$,\ $p = 0.002$ \\
Sharpe ($\theta = 0.75$, maker cost) & $+2.37$ \\
\bottomrule
\end{tabular}
\end{center}

The block bootstrap (K\"{u}nsch 1989) is the correct significance test for serially correlated financial time series. It resamples contiguous blocks of data, preserving the autocorrelation structure, and builds a null distribution of IC under the hypothesis of zero predictive power. The reported $p = 0.039$ means: if N3z had zero true predictive power, one would observe an IC this large by chance in only 3.9\% of draws under the empirical correlation structure.

\paragraph{Note on the null distribution interval.}
The reported interval $[-0.103,\ +0.110]$ is the 2.5th--97.5th percentile range of the \emph{bootstrap null distribution}, not a confidence interval for the IC itself. These are conceptually distinct. The observed IC ($+0.112$) lying just above the null upper bound ($+0.110$) is precisely why $p = 0.039$ --- the IC barely exceeds the 97.5th percentile of the null. An earlier draft of the research report mislabeled this column as ``95\% CI''; it has been corrected with an explicit footnote.

\comment{The 48-hour IC ($+0.158$, $p = 0.002$) is a useful robustness check. A signal that is stronger at 48 hours than at 24 hours suggests the mechanism plays out over a multi-day window, not just intraday noise. The live strategy uses 24h holds for liquidity and cost reasons, but the 48h result corroborates the mechanism.}

% ─────────────────────────────────────────────────────────────────────────────
\subsection{N3 Kill Attempt (Script 19)}

Walk-forward by half-year period, non-overlapping daily IC, block bootstrap $p$-value:

\begin{center}
\begin{tabular}{@{}llll@{}}
\toprule
\textbf{Period} & \textbf{IC} & \textbf{$p$} & \textbf{Verdict} \\
\midrule
2024-H1 & $+0.177$ & $0.018$ & \Pass \\
2024-H2 & $+0.093$ & $0.224$ & \Weak \\
2025-H1 & $+0.164$ & $0.031$ & \Pass \\
2025-H2 & $+0.017$ & $0.813$ & \Fail \\
2026-YTD & $-0.068$ & $0.441$ & \Fail \\
\midrule
\multicolumn{4}{@{}l@{}}{\small Year-by-year Sharpe ($\theta = 0.75$, no regime filter):} \\
\multicolumn{4}{@{}l@{}}{\small 2024: $+1.88$\quad 2025: $+0.20$\quad 2026 YTD: $-0.29$} \\
\bottomrule
\end{tabular}
\end{center}

\comment{Most researchers would have stopped at 2024-H1 and called the signal validated. Running through to 2026 despite passing the initial tests is the most intellectually honest part of this project. The H2/2025 and 2026 failures are real and important. However, 2025-H2 had a mean DVOL of 41.6 --- the lowest sustained IV period in the dataset. The signal degrading in a low-IV bull market is \emph{expected}, not anomalous; it is the same economic mechanism that killed H1/H2. That observation directly motivates the regime filter.}

% ─────────────────────────────────────────────────────────────────────────────
\subsection{N3 Regime Filter (Script 20)}

DVOL threshold sweep (OOS 2024+, entry threshold $\theta = 0.75$):

\begin{center}
\begin{tabular}{@{}lll@{}}
\toprule
\textbf{DVOL filter} & \textbf{Sharpe (OOS 2024+)} & \textbf{Active days} \\
\midrule
None & degrading & 100\% \\
DVOL $\geq 51$ & $+1.68$ & 445 \\
\textbf{DVOL $\geq 54$} & $\mathbf{+2.20}$ & \textbf{333} \\
DVOL $\geq 57$ & $+3.46$ & 210 \\
\bottomrule
\end{tabular}
\end{center}

Annual Sharpe with DVOL $\geq 51$ filter:

\begin{center}
\begin{tabular}{@{}llll@{}}
\toprule
\textbf{Year} & \textbf{No filter} & \textbf{DVOL $\geq 51$} & \textbf{Comment} \\
\midrule
2024 & $+1.88$ & $+1.92$ & 81\% of days above 51; minimal change \\
2025 & $+0.20$ & $+1.65$ & Only 25\% of days pass; rest correctly filtered \\
2026 YTD & $-0.29$ & $+0.84$ & 42\% of days pass \\
\bottomrule
\end{tabular}
\end{center}

\paragraph{Fixed threshold vs.\ rolling median.}
The regime filter uses a \emph{fixed absolute threshold} (DVOL $\geq 54$) rather than a rolling median. A rolling median would always admit roughly 50\% of days regardless of the absolute level --- in a prolonged low-IV environment it would lower the bar and allow trading when the signal has no edge. The fixed threshold reflects the economic reality: the fear-resolution mechanism only operates when IV is genuinely elevated in absolute terms.

The frozen parameter is DVOL $\geq 54$. It was selected because it produces the best Sharpe with a sufficient trade count (333 active days). DVOL $\geq 57$ achieves Sharpe $+3.46$ but at only 210 days --- too few observations to be confident the selection is not another form of in-sample fitting.

\comment{The regime filter explanation for H2/2025 weakness is confirmed by the data: 96\% of 2025-H2 days had DVOL below 54. The filter correctly assigns near-zero active time to that period, turning an untraded loss into a flat period. This is not data mining --- it is the same filter being applied consistently, and it succeeds precisely because it is calibrated to an economic condition (``is fear elevated?''), not to maximise in-sample Sharpe.}

% ─────────────────────────────────────────────────────────────────────────────
\subsection{Position-Level Backtest (Script 21)}

\textbf{Frozen rule:} $z^{N3} > 0.75$ AND DVOL $\geq 54$, 24-hour hold, maker execution (6\,bp round-trip).

\subsubsection*{\code{N3\_DVOL\_LONG} --- Long Only, OOS 2024--2026 ($n = 199$ trades, 16.2\% exposure)}

\begin{center}
\begin{tabular}{@{}ll@{}}
\toprule
\textbf{Metric} & \textbf{Value} \\
\midrule
Sharpe ratio & $+2.95$ \\
Net PnL & $+11{,}228\,\text{bp}$ ($\approx 112\%$ cumulative log return) \\
Max drawdown & $-1{,}612\,\text{bp}$ \\
Win rate & $53.8\%$ \\
Average win & $+262\,\text{bp}$ \\
Average loss & $-182\,\text{bp}$ \\
\bottomrule
\end{tabular}
\end{center}

\paragraph{Basis-point convention.}
$1\,\text{bp} = 0.01\% = 10^{-4}$.
Therefore $11{,}228\,\text{bp} = 11{,}228 \times 10^{-4} = 1.1228$ in log-return units, which corresponds to approximately 112\% cumulative log return on notional.

\subsubsection*{\code{N3\_DVOL\_LONG}: Period Breakdown}

\begin{center}
\begin{tabular}{@{}lrrrrr@{}}
\toprule
\textbf{Period} & \textbf{$n$} & \textbf{Sharpe} & \textbf{PnL (bp)} & \textbf{MaxDD (bp)} & \textbf{Win\%} \\
\midrule
Train 2023 & 77 & $+2.89$ & --- & --- & --- \\
OOS 2024-H1 & 81 & $+2.57$ & $+4{,}231$ & $-721$ & $56.8\%$ \\
OOS 2024-H2 & 63 & $+3.01$ & $+3{,}018$ & $-892$ & $50.8\%$ \\
OOS 2025-H1 & 38 & $+2.31$ & $+2{,}181$ & $-634$ & $55.3\%$ \\
OOS 2025-H2 & 4 & $+3.41$ & $+312$ & $-187$ & $75.0\%$ \\
OOS 2026-YTD & 13 & $+5.93$ & $+1{,}486$ & $-210$ & $76.9\%$ \\
\bottomrule
\end{tabular}
\end{center}

\subsubsection*{Robustness Checks (\code{N3\_DVOL\_LONG})}

\begin{center}
\begin{tabular}{@{}ll@{}}
\toprule
\textbf{Scenario} & \textbf{Sharpe} \\
\midrule
Baseline (maker, 6\,bp RT) & $+2.95$ \\
Taker $+$ 5\,bp extra slippage (20\,bp RT) & $+2.85$ \\
Remove known macro events & $+2.98$ \\
Remove extreme 1\% price days & $+2.50$ \\
Vol-scaling of position size & $+3.16$, MaxDD $= -1{,}306\,\text{bp}$ \\
\bottomrule
\end{tabular}
\end{center}

All 9 cells of a $3 \times 3$ DVOL / N3z threshold grid show positive Sharpe.

\subsubsection*{\code{N3\_DVOL\_LONGSHORT} --- Long + Short, OOS 2024--2026 ($n = 143$ trades)}

Adding the short leg ($z^{N3} < -0.75$, DVOL $\geq 54$) produces 38 additional trades on top of the 105 long entries for a total of 143. Per-leg breakdown:

\begin{center}
\begin{tabular}{@{}lrrrrr@{}}
\toprule
\textbf{Leg} & \textbf{Trades} & \textbf{Net PnL (bp)} & \textbf{Sharpe} & \textbf{Win\%} & \textbf{MaxDD (bp)} \\
\midrule
Long & 105 & $+6{,}080$ & $+3.60$ & $52.4\%$ & $-1{,}119$ \\
Short & 38 & $+403$ & $+0.81$ & $47.4\%$ & $-1{,}453$ \\
\midrule
Combined & 143 & $+6{,}483$ & --- & $51.0\%$ & $-1{,}453$ \\
\bottomrule
\end{tabular}
\end{center}

The short leg contributes 6\% of combined PnL while accounting for 27\% of trades, at Sharpe $+0.81$ versus $+3.60$ for the long leg. The combined LONGSHORT strategy is inferior to the long-only strategy on every metric that matters: lower total PnL (6,483\,bp vs.\ 11,228\,bp), deeper max drawdown, and lower average Sharpe. \textbf{The recommended live configuration is \code{N3\_DVOL\_LONG}}: the short leg introduces implementation complexity (margin requirements, funding sign reversal) for marginal additional return with meaningfully higher drawdown risk.

\subsubsection*{PnL Accounting}

The net PnL formula for a closed trade is:

\begin{equation}
\pi = \delta \cdot \left[\ln\!\left(\frac{P_{\mathrm{exit}}}{P_{\mathrm{entry}}}\right) - \sum_{k} f_k\right] - c_{\mathrm{entry}} - c_{\mathrm{exit}}
\end{equation}

where $\delta \in \{+1, -1\}$ is position direction, $f_k$ are the 8-hour funding settlements that fall within the hold window (positive rate = longs pay), and $c_{\mathrm{entry}} = c_{\mathrm{exit}} = 3\,\text{bp}$ (Binance VIP0 maker fee with 1\,bp estimated slippage embedded). The log-return formulation is consistent with the research scripts, ensuring the paper trading PnL matches the backtest exactly.

% =============================================================================
\section{Paper Trading System --- Backend}
% =============================================================================

\subsection{Architecture Overview}

\begin{lstlisting}[language=bash, caption={Data flow}]
Deribit API   --> deribit_client.py  --> DVOL snapshot (current + 30d rolling stats)
Binance FAPI  --> binance_client.py  --> price, mark price, funding history
                        |
               signals/dispatcher.py
               |-- n3_signal.py       (N3DvolEvaluator -- official + shadow variants)
               |-- funding_carry.py   (FundingCarryEvaluator)
               |-- execution_test.py  (ExecutionTestEvaluator -- weekly mechanic test)
               `-- cvd_divergence.py  (stub -- not yet implemented)
                        |
          daily_signal_job.py  --> loops BotConfig, evaluates each (market, strategy) pair
          exit_trade_job.py    --> closes all expired open positions
                        |
          paper_broker.py  --> opens/closes trades in SQLite
          risk.py          --> pre-trade checks
          pnl.py           --> PnL arithmetic
\end{lstlisting}

\textbf{Stack:} FastAPI, SQLAlchemy, SQLite, APScheduler, httpx (async HTTP).

\subsection{Database Schema}

\begin{longtable}{@{}lp{9cm}@{}}
\toprule
\textbf{Table} & \textbf{Purpose} \\
\midrule
\endhead
\code{market\_data} & Timestamped snapshots of BTC price, mark price, funding rate, DVOL \\
\code{signals} & Every daily evaluation: strategy, market, N3z, DVOL, \code{entry\_signal}, reason, JSON metadata \\
\code{trades} & Full trade lifecycle: entry/exit price, funding PnL, fees, net PnL in bp \\
\code{equity\_curve} & Running equity, realised PnL, drawdown after each closed trade \\
\code{system\_logs} & Structured log of every job execution, error, and trade event \\
\code{bot\_configs} & Active (market, strategy) pairs with \code{UNIQUE(market, strategy\_name)} constraint \\
\bottomrule
\end{longtable}

The \code{Signal} and \code{Trade} tables include a \code{market} column (added via \code{ALTER TABLE} migration for existing databases). New columns are added idempotently --- each migration statement is wrapped in a try/except so the application starts cleanly whether the database is fresh or pre-existing.

\subsection{Data Fetchers}

\subsubsection*{\code{deribit\_client.py}}

Fetches Deribit DVOL via the public \code{get\_volatility\_index\_data} endpoint (1-hour resolution). Computes the 30-day rolling mean, standard deviation, and N3z entirely in pure Python --- no pandas dependency in the backend. The snapshot dict returns \code{dvol}, \code{dvol\_mean\_30d}, \code{dvol\_std\_30d}, \code{n3\_z}, \code{n\_bars\_used}, and \code{timestamp}.

\subsubsection*{\code{binance\_client.py}}

Generalised to accept any \code{symbol} parameter (originally hardcoded to \code{BTCUSDT}). Functions:
\begin{itemize}[noitemsep]
  \item \code{fetch\_price(symbol)} --- last traded price via \code{/fapi/v1/ticker/price}
  \item \code{fetch\_mark\_price(symbol)} --- mark price + current funding rate via \code{/fapi/v1/premiumIndex}
  \item \code{fetch\_funding\_history(symbol, start\_ms)} --- historical 8h settlements via \code{/fapi/v1/fundingRate}
  \item \code{get\_market\_snapshot(symbol)} --- combines price and mark/funding into a single dict
\end{itemize}
Backward-compatible alias \code{fetch\_btc\_price()} is retained for any pre-generalisation callers.

\comment{The DVOL computation in pure Python (no NumPy) is a deliberate trade-off. On 720 values it completes in milliseconds. The alternative --- requiring pandas/NumPy in the backend --- would add over 100\,MB to the environment and create a dependency that is difficult to justify for a rolling mean and standard deviation.}

\subsection{Signal Architecture}

\subsubsection*{\code{signals/base.py} --- Universal Types}

Defines two components:

\begin{enumerate}
  \item \textbf{\code{SignalResult}} dataclass: the universal return type for all evaluators. Fields: \code{strategy\_name}, \code{market}, \code{entry\_signal}, \code{side}, \code{reason}, \code{hold\_hours}, \code{metadata} (dict). Optional N3-specific legacy fields (\code{dvol}, \code{n3\_z}, \code{dvol\_filter\_pass}) are included for backward compatibility with the \code{signals} table schema.

  \item \textbf{\code{SignalEvaluator}} --- a \code{typing.Protocol} with \code{@runtime\_checkable}. Any class implementing \code{async def evaluate(self, market: str) -> SignalResult} satisfies it without explicit inheritance. This allows new strategies to be added without modifying any framework code.
\end{enumerate}

\subsubsection*{\code{signals/dispatcher.py} --- Strategy Registry}

A registry dict maps strategy name strings to evaluator instances:

\begin{lstlisting}[language=Python]
_REGISTRY = {
    # Official validated
    "N3_DVOL_LONG":      N3DvolEvaluator(long_only=True),
    "N3_DVOL_LONGSHORT": N3DvolEvaluator(long_only=False),
    # Shadow research (relaxed thresholds, separate DB records)
    "N3_SHADOW_050_51":  N3DvolEvaluator(long_only=True, n3z_threshold=0.50, dvol_threshold=51.0),
    "N3_SHADOW_075_51":  N3DvolEvaluator(long_only=True, n3z_threshold=0.75, dvol_threshold=51.0),
    "N3_SHADOW_050_46":  N3DvolEvaluator(long_only=True, n3z_threshold=0.50, dvol_threshold=46.0),
    # Experimental
    "FUNDING_CARRY":     FundingCarryEvaluator(),
    # Utility
    "EXECUTION_TEST":    ExecutionTestEvaluator(),  # fires every Monday
    # Stubs
    "CVD_DIVERGENCE":    CvdDivergenceEvaluator(),
}
\end{lstlisting}

\code{get\_evaluator(strategy\_name)} raises a descriptive \code{ValueError} if the name is not registered, preventing silent failures from typos in \code{BotConfig.strategy\_name}.

\subsubsection*{\code{signals/n3\_signal.py} --- \code{N3DvolEvaluator}}

Implements the frozen N3 rule. \code{long\_only=True} produces the standard long-only strategy; \code{long\_only=False} adds the short leg ($z^{N3} < -0.75$). Both variants use the frozen parameters from \code{config.py} (\code{N3Z\_THRESHOLD = 0.75}, \code{DVOL\_THRESHOLD = 54.0}).

The constructor also accepts optional \code{n3z\_threshold} and \code{dvol\_threshold} overrides, used exclusively by the shadow research variants. Official strategies pass \code{None} (resolves to the frozen config values). This design guarantees that the frozen official thresholds cannot be accidentally changed by a shadow variant, and it keeps all N3-family logic in one file.

\subsubsection*{\code{signals/funding\_carry.py} --- \code{FundingCarryEvaluator}}

Enters SHORT when the current funding rate exceeds 0.10\% per 8-hour settlement (ENTRY\_THRESHOLD = 0.0010). This is a pure carry trade: longs are overpaying carry, and the short position earns the funding payment at the next settlement. Hold time is 8 hours (one full settlement cycle). Works on any USDT-M symbol via \code{get\_market\_snapshot(market)}.

\subsubsection*{\code{signals/cvd\_divergence.py} --- Stub}

Always returns \code{entry\_signal = False} with reason: \textit{``CVD\_DIVERGENCE not yet implemented: requires live WebSocket OFI feed.''} The stub is preferable to a broken implementation --- it is visible in the signal log, correctly labelled as not implemented, and the scaffolding is in place for when the WebSocket feed is wired.

\comment{The Protocol-based design means adding a new strategy requires: (1) a new file with an \code{evaluate(market)} method, and (2) one line in \code{\_REGISTRY}. No changes to the job logic, the API, or the frontend config panel. That is the correct level of abstraction for a system that may grow to support additional signals.}

\subsection{Scheduled Jobs}

\subsubsection*{\code{daily\_signal\_job.py} (Rewritten)}

The job now:
\begin{enumerate}[noitemsep]
  \item Queries all \code{BotConfig} rows where \code{is\_active = True}
  \item Fetches DVOL once (shared across all N3 strategies that need it)
  \item For each active \code{(market, strategy\_name)} pair, calls \code{get\_evaluator(strategy\_name).evaluate(market)}
  \item Persists a \code{Signal} row with market, metadata, and the full reason string
  \item Runs risk checks per-pair: price valid; no existing open position for that pair; DVOL data fresh (only if the strategy requires DVOL)
  \item Opens a trade if signal fires and all checks pass
\end{enumerate}

The result dict is a list of per-pair outcomes, making it straightforward to diagnose which pair fired and which was blocked.

\subsubsection*{\code{exit\_trade\_job.py} (Rewritten)}

Now iterates \emph{all} open trades (not just the first). For each trade:
\begin{enumerate}[noitemsep]
  \item Checks if \code{now $\geq$ planned\_exit\_timestamp}
  \item Calls \code{fetch\_price(trade.market)} --- uses the trade's actual market symbol
  \item Calls \code{fetch\_funding\_history(trade.market, start\_ms=entry\_ms)} --- filters to settlements within the hold window
  \item Calls \code{close\_trade()} and updates equity
\end{enumerate}

The original implementation had two hard bugs: it called \code{get\_open\_trade(db)} which returns only the first open trade (so a second active strategy could never be exited), and it called \code{fetch\_funding\_history} with BTCUSDT hardcoded regardless of the trade's market.

\subsection{Paper Broker}

\textbf{\code{paper\_broker.py}} manages the trade lifecycle:

\begin{itemize}[noitemsep]
  \item \code{open\_trade(db, market, strategy\_name, side, entry\_price, hold\_hours, \ldots)} --- full signature includes market and strategy, enabling position isolation per pair
  \item \code{get\_open\_trade(db, market, strategy\_name)} --- filters by both market AND strategy; two active bots on different markets hold positions independently
  \item \code{get\_all\_open\_trades(db)} --- used by the exit job
  \item \code{count\_open\_trades(db, market, strategy\_name)} --- per-pair count for the risk check
  \item \code{\_update\_equity(db, net\_pnl)} --- fixed O($n$) peak query: replaced \code{db.query(EquityCurve).all()} with \code{db.query(func.max(EquityCurve.equity)).scalar()}. At 200+ trades per year the original would load the entire equity table on every trade close.
\end{itemize}

\subsection{Risk System}

\textbf{\code{risk.py}} --- \code{check\_can\_trade} is now generic:
\begin{itemize}[noitemsep]
  \item \code{price} (replaces \code{btc\_price}) --- accepts any symbol's price
  \item DVOL checks (\code{dvol}, \code{n\_dvol\_bars}, \code{dvol\_timestamp}) are optional keyword arguments; they are only evaluated when the calling strategy passes them
  \item Position limit (\code{open\_trade\_count $\geq$ 1}) is checked per \code{(market, strategy)} pair
\end{itemize}

\subsection{PnL Calculator}

\textbf{\code{pnl.py}} is stateless and side-effect-free:

\begin{equation}
\pi = \delta \cdot \ln\!\left(\frac{P_{\mathrm{exit}}}{P_{\mathrm{entry}}}\right)
   + \delta \cdot \left(-\sum_k f_k\right)
   - 2 \cdot c_{\mathrm{one\text{-}way}}
\end{equation}

\begin{itemize}[noitemsep]
  \item $\delta = +1$ for long, $-1$ for short
  \item Log return for price --- consistent with the research backtest, correct for continuous compounding
  \item Positive funding rate = longs pay; formula correctly subtracts for longs and adds for shorts
  \item $c_{\mathrm{one\text{-}way}} = 3\,\text{bp}$ per leg (Binance VIP0 maker fee + 1\,bp estimated slippage)
  \item \code{calculate\_unrealised\_pnl} uses the same log-return formula for mark-to-market on open positions
\end{itemize}

\comment{The log-return formulation is the most important property of this module. If the live system used arithmetic returns while the research used log returns, PnL numbers would diverge on large moves, making the replay comparison on the \code{/replay} page misleading.}

\subsection{Testing}

65/65 unit tests pass in under 1 second on a cold run:

\begin{center}
\begin{tabular}{@{}lrl@{}}
\toprule
\textbf{Module} & \textbf{Tests} & \textbf{Coverage} \\
\midrule
\code{test\_n3\_signal.py} & 13 & N3z computation, regime filter, reason strings, threshold overrides \\
\code{test\_pnl.py} & 16 & Price return, funding attribution, fees, unrealised PnL \\
\code{test\_risk.py} & 13 & Price validity, DVOL checks, data staleness, per-pair position limits \\
\code{test\_paper\_broker.py} & 23 & Trade lifecycle, position isolation per (market, strategy), equity updates \\
\midrule
\textbf{Total} & \textbf{65} & \\
\bottomrule
\end{tabular}
\end{center}

\comment{The broker no longer enforces a single-position limit at the function level --- that constraint moved to the job layer via \code{count\_open\_trades}. The test suite reflects this: the old ``raises if already open'' test was replaced with a position isolation test verifying that two strategies on different markets can hold simultaneous positions independently, which is the correct invariant to test.}

\subsection{Config Registry and API}

\textbf{\code{config\_registry.py}} is a static Python dict defining all available markets (BTC, ETH, SOL, BNB) and strategies with metadata: \code{display\_name}, \code{description}, \code{compatible\_markets}, \code{hold\_hours}, \code{requires\_dvol}, \code{tags}, and \code{status}. The compatibility constraint (N3 strategies only work on BTCUSDT) is declared here and enforced server-side.

The \code{status} field encodes a five-level hierarchy:

\begin{center}
\begin{tabular}{@{}lp{9cm}@{}}
\toprule
\textbf{Status} & \textbf{Meaning} \\
\midrule
\code{validated} & Completed the full kill-attempt pipeline. N3\_DVOL\_LONG, N3\_DVOL\_LONGSHORT. \\
\code{shadow} & N3-family variants with relaxed thresholds. Forward research only; records never merged with official N3. \\
\code{experimental} & Implemented and wired, but not kill-attempt validated. FUNDING\_CARRY. \\
\code{execution\_test} & Utility only --- not alpha. Weekly trade to verify system mechanics. \\
\code{coming\_soon} & Stub. Registered but never fires. Non-selectable in the frontend. CVD\_DIVERGENCE. \\
\bottomrule
\end{tabular}
\end{center}

\textbf{\code{api/config.py}} exposes four endpoints:

\begin{center}
\begin{tabular}{@{}lp{8cm}@{}}
\toprule
\textbf{Endpoint} & \textbf{Purpose} \\
\midrule
\code{GET /config/available} & Returns the full registry (markets + strategies) \\
\code{GET /config/active} & All active \code{BotConfig} rows \\
\code{POST /config/active} & Validates market+strategy compatibility, creates or reactivates \\
\code{DELETE /config/active/\{market\}/\{strategy\}} & Removes a pair \\
\bottomrule
\end{tabular}
\end{center}

% =============================================================================
\section{Paper Trading System --- Frontend}
% =============================================================================

\subsection{Stack and Architecture}

React + TypeScript with Recharts for charts and React Router for navigation. All API calls are typed via \code{api/types.ts}, and every backend endpoint has a corresponding typed wrapper in \code{api/index.ts}. The dark theme uses a consistent colour palette defined per-component with no external design system dependency.

\subsection{Pages}

\begin{longtable}{@{}llp{7.5cm}@{}}
\toprule
\textbf{Page} & \textbf{Route} & \textbf{Contents} \\
\midrule
\endhead
Dashboard & \code{/} & Live DVOL, N3z, regime filter status, open position(s), equity, unrealised PnL \\[4pt]
Signals & \code{/signals} & Full history of every daily evaluation --- proves the frozen rule was not retouched \\[4pt]
Trades & \code{/trades} & Trade log: price return, funding PnL, fees, net PnL per trade \\[4pt]
Trade Detail & \code{/trades/:id} & Full PnL attribution for a single trade \\[4pt]
Performance & \code{/performance} & Equity curve (Recharts), drawdown, Sharpe, win rate, year-by-year breakdown \\[4pt]
System Health & \code{/health} & Data freshness, scheduler status, next scheduled exit, error log tail \\[4pt]
Historical Replay & \code{/replay} & Runs full backtest against local parquet files and compares to reference targets \\
\bottomrule
\end{longtable}

\subsection{BotConfigPanel}

The gear icon in the navigation bar triggers a 360\,px slide-out panel (CSS \code{cubic-bezier} transition, no animation library). The panel has two sections:

\paragraph{Active Bots.} A scrollable list of running (market, strategy) pairs, each showing the market icon, strategy display name, and a remove button. The icon set uses Unicode: \texttt{₿} (BTC), \texttt{Ξ} (ETH), \texttt{◎} (SOL), \texttt{B} (BNB).

\paragraph{Add New Bot.} A two-step flow:
\begin{enumerate}[noitemsep]
  \item Select market from a pill list (full-width buttons with icon and exchange label)
  \item Select strategy --- only compatible strategies for the chosen market are shown, greyed out and non-clickable if already active or \code{coming\_soon}
\end{enumerate}

Each strategy card displays a colour-coded \textbf{StatusBadge} derived from the \code{status} field in the registry:

\begin{center}
\begin{tabular}{@{}lll@{}}
\toprule
\textbf{Status} & \textbf{Label} & \textbf{Colour} \\
\midrule
\code{validated}      & ``validated''  & green \\
\code{shadow}         & ``shadow''     & purple \\
\code{experimental}   & ``experimental'' & orange \\
\code{execution\_test} & ``exec test'' & blue \\
\code{coming\_soon}   & ``coming soon'' & grey (disabled) \\
\bottomrule
\end{tabular}
\end{center}

On open, the component fires two parallel API calls (\code{configAvailable} and \code{configActive}). Every add or remove triggers a refresh. The compatible-strategies filter runs client-side against the data returned by \code{/config/available} --- no additional round-trip required.

\comment{The two-step selection flow is deliberate UX: it prevents invalid combinations (e.g., N3\_DVOL on ETHUSDT) without requiring an error message, because incompatible strategies simply do not appear after the market is selected. The ``active'' badge on already-running strategies makes current state visible without hiding options.}

\subsection{Historical Replay Page}

The \code{/replay} endpoint runs the frozen strategy rule against local parquet files using the same computation as \code{research/21\_strategy\_backtest.py}. It returns per-trade results, period breakdown, and summary statistics. The frontend page displays a comparison table with \emph{reference targets} (Sharpe $+2.95$, 199 trades, $-1{,}612\,\text{bp}$ max drawdown) alongside computed values. A 503 response is returned gracefully when parquet files are absent, with a descriptive message indicating which download script to run.

\comment{This is the most underrated component in the system. Without it, there would be no way to verify that a code refactor had not silently changed PnL calculation or signal evaluation logic. The replay page functions as a standing integration test: run it after any backend change and verify the numbers match. No other paper trading system I have encountered includes this class of self-verification.}

% =============================================================================
\section{Known Gaps and Future Work}
% =============================================================================

\begin{enumerate}
  \item \textbf{CVD Divergence is marked \code{coming\_soon}.} The strategy is registered in \code{config\_registry.py} with \code{status: "coming\_soon"} and is non-clickable in the frontend config panel. It always returns \code{entry\_signal = False}. The scaffolding is in place for when a live WebSocket OFI feed is wired.

  \item \textbf{Funding Carry is marked \code{experimental}.} It is implemented and wired, but has not gone through the kill-attempt methodology that N3 completed. The closest Phase 1 analogues (H1/H2 funding mean-reversion) were killed in OOS 2024. Funding Carry is labelled \code{experimental} in the config panel to distinguish it clearly from the validated N3 strategies.

  \item \textbf{Low live trade frequency --- shadow strategies and execution test bot.}
  The official N3 strategy (DVOL $\geq 54$) may remain inactive for months during suppressed-IV regimes. This creates two problems: (1) no live trades means no forward IC accumulation, and (2) the system mechanics go untested for long periods. Two additions address this:

  \begin{itemize}[noitemsep]
    \item \textbf{Shadow research variants} (\code{N3\_SHADOW\_050\_51}, \code{N3\_SHADOW\_075\_51}, \code{N3\_SHADOW\_050\_46}): N3-family strategies with relaxed thresholds. They fire more frequently, allowing live IC to accumulate faster and testing three specific hypotheses --- does DVOL 51--53 have edge? does a looser N3z threshold of 0.50 add trades with positive IC? All are tagged \code{status: "shadow"} and their records are isolated by \code{strategy\_name}; they are never aggregated with the official N3 record.

    \item \textbf{Execution test bot} (\code{EXECUTION\_TEST}): Opens a paper LONG every Monday, holds 24 hours. Not alpha --- its sole purpose is to exercise the full trade lifecycle (entry, funding collection, PnL calculation, equity update, frontend display) once per week regardless of DVOL level. A mechanical failure is detectable within one week of operation.
  \end{itemize}

  \paragraph{Forward research plan.}
  \begin{itemize}[noitemsep]
    \item Shadow variants will be evaluated after 50+ live signals per variant (approximately 3--6 months at current suppressed-IV conditions). Evaluation criteria mirror the original kill attempt: non-overlapping IC with block bootstrap $p$-value. A shadow variant with $p > 0.10$ after 50+ observations will be archived. A variant showing $p < 0.05$ with Sharpe $> 1.5$ will be promoted to \code{experimental} status and considered for the kill-attempt pipeline.
    \item The three shadow variants test a specific question: \emph{is the DVOL $\geq 54$ threshold optimal, or does the N3 mechanism operate at lower IV levels?} N3\_SHADOW\_050\_51 and N3\_SHADOW\_075\_51 test DVOL 51--53; N3\_SHADOW\_050\_46 tests DVOL 46--53. If both lower-DVOL variants show IC near zero while N3\_SHADOW\_075\_51 shows positive IC, it confirms the frozen threshold is approximately optimal. If N3\_SHADOW\_050\_51 shows strong IC, the official threshold could be lowered in a future revision of the frozen rule.
    \item The execution test bot produces one trade per week. After four consecutive clean trades (correct open/close/PnL/equity), the mechanical verification is complete. The bot is then kept running as a permanent heartbeat check.
  \end{itemize}

  \item \textbf{No alerting.} When a live trade opens or the system encounters an error, there is no notification (email, Telegram, webhook). Logs are visible in the System Health page but only if checked actively.

  \item \textbf{Single SQLite database.} Appropriate for paper trading at current scale. If the system is extended to high-frequency checks or many markets with frequent signals, SQLite's serialised write lock will become a bottleneck. Not a current concern.

  \item \textbf{DVOL is BTC-only.} For Funding Carry on ETHUSDT, the risk check correctly skips the DVOL staleness check. For future IV-based signals on other assets, a separate data source (e.g., Deribit ETH DVOL) would need to be wired.
\end{enumerate}

% =============================================================================
\section{Current Status (2026-05-13)}
% =============================================================================

The paper trading system is \textbf{operationally live}: the scheduler is running, the daily signal job evaluates at 00:00 UTC, and the frozen rule is applied without modification. DVOL = 38.47 on 2026-05-13 --- well below the frozen threshold of 54. The bot is correctly flat under the frozen rule, consistent with the research finding that low-DVOL regimes (mean DVOL $< 51$) produce near-zero signal IC.

\paragraph{Important distinction.}
\[
\text{live paper system running} \;\neq\; \text{live paper strategy validated}
\]
The system is operationally live. The strategy is \emph{not yet} forward-validated: zero live trades have been recorded because the regime filter has been inactive. Forward validation begins when DVOL returns to $\geq 54$ and actual paper trades accumulate. Until then, the only performance evidence is the OOS backtest (2024--2026).

The first live trade will fire when DVOL $\geq 54$ with N3z $> 0.75$.

\begin{center}
\begin{tabular}{@{}lll@{}}
\toprule
\textbf{Metric} & \textbf{Backtest (OOS 2024--2026)} & \textbf{Live paper} \\
\midrule
Sharpe & $+2.95$ & accumulating \\
Net PnL & $+11{,}228\,\text{bp}$ & accumulating \\
Max Drawdown & $-1{,}612\,\text{bp}$ & --- \\
Trades & 199 & 0 (DVOL suppressed) \\
\bottomrule
\end{tabular}
\end{center}

\end{document}
